\chapter{User Documentation}
\label{ch:user}

This chapter provides comprehensive documentation for end users of the Urban Air Quality Forecasting application. The application is built as a multi-page Streamlit web application that enables users to explore air quality data, generate forecasts using machine learning models, and monitor urban air pollution trends through interactive dashboards.

\section{System Requirements}
\label{sec:user-requirements}

\subsection{Hardware Requirements}

\begin{table}[htbp]
\centering
\caption{Minimum and recommended hardware specifications}
\label{tab:hardware-requirements}
\begin{tabular}{|l|l|l|}
\hline
\textbf{Component} & \textbf{Minimum} & \textbf{Recommended} \\
\hline
Processor & Dual-core 2.0 GHz & Quad-core 3.0 GHz+ \\
\hline
RAM & 4 GB & 8 GB+ \\
\hline
Storage & 2 GB free space & 5 GB+ free space \\
\hline
Network & Broadband connection & Stable broadband (10 Mbps+) \\
\hline
GPU & Not required & CUDA-compatible (for training) \\
\hline
\end{tabular}
\end{table}

\subsection{Software Requirements}

\begin{itemize}
    \item \textbf{Operating System:} Windows 10/11, macOS 12+, or Ubuntu 20.04+
    \item \textbf{Python:} Version 3.9 or higher (3.11 recommended)
    \item \textbf{pip:} Version 21.0 or higher
    \item \textbf{Git:} Version 2.30+ (for cloning the repository)
    \item \textbf{Database:} SQLite 3.35+ (bundled with Python) or PostgreSQL 13+ (optional, for production)
\end{itemize}

\subsection{Browser Requirements}

The application is tested and supported on the following browsers:

\begin{itemize}
    \item Google Chrome 90+ (recommended)
    \item Mozilla Firefox 88+
    \item Microsoft Edge 90+
    \item Safari 14+
\end{itemize}

\noindent JavaScript must be enabled, and a minimum viewport width of 1024 pixels is recommended for optimal dashboard rendering. The application is responsive but designed primarily for desktop use.

\section{Installation and Configuration}
\label{sec:user-installation}

\subsection{Cloning the Repository}

Begin by cloning the project repository to your local machine:


\subsection{Setting Up the Virtual Environment}

It is strongly recommended to use a Python virtual environment to isolate project dependencies from the system Python installation.

\subsubsection{On Linux/macOS}


\subsubsection{On Windows}


\subsection{Installing Dependencies}

The project uses a \texttt{requirements.txt} file to manage all Python dependencies:


\noindent The key dependencies include:

\begin{table}[htbp]
\centering
\caption{Core project dependencies}
\label{tab:dependencies}
\begin{tabular}{|l|l|l|}
\hline
\textbf{Package} & \textbf{Version} & \textbf{Purpose} \\
\hline
streamlit & $\geq$ 1.28.0 & Web application framework \\
\hline
pandas & $\geq$ 2.0.0 & Data manipulation \\
\hline
scikit-learn & $\geq$ 1.3.0 & Machine learning models \\
\hline
xgboost & $\geq$ 2.0.0 & Gradient boosting models \\
\hline
plotly & $\geq$ 5.17.0 & Interactive visualizations \\
\hline
google-generativeai & $\geq$ 0.3.0 & Gemini API integration \\
\hline
bcrypt & $\geq$ 4.0.0 & Password hashing \\
\hline
streamlit-option-menu & $\geq$ 0.3.6 & Navigation components \\
\hline
\end{tabular}
\end{table}

\subsection{Configuring Secrets}

Streamlit uses a secrets management system via the \texttt{.streamlit/secrets.toml} file. Create this file in the project root:


\noindent Populate the secrets file with the required configuration:


\noindent\textbf{Important:} Never commit the \texttt{secrets.toml} file to version control. Ensure it is listed in your \texttt{.gitignore} file.

\subsection{Database Initialization}

The application automatically initializes the SQLite database on first run. However, you can manually trigger initialization:


\noindent This creates the necessary tables for user accounts, prediction history, and application settings.

\subsection{Model Training}

Before generating predictions, the machine learning models must be trained on historical air quality data. The training pipeline supports multiple model architectures:


\noindent Upon successful training, model artifacts are saved to the \texttt{models/} directory as serialized \texttt{.pkl} files. Training metrics (RMSE, MAE, R\textsuperscript{2}) are logged to the console and saved to \texttt{models/training\_report.json}.

\subsection{Running the Application}

Launch the application using the Streamlit CLI:

\begin{verbatim}
streamlit run front_end/Home.py
\end{verbatim}

\noindent The application will be accessible at \texttt{http://localhost:8501}. On first launch, Streamlit may ask permission to collect usage statistics---this is optional and can be declined.

\subsection{Production Deployment Considerations}

For production environments, additional configuration is recommended:



\section{Authentication System}
\label{sec:user-auth}

The application implements a comprehensive authentication system to protect user data and personalize the forecasting experience. All passwords are hashed using the bcrypt algorithm before storage, ensuring that plaintext credentials are never persisted.

\subsection{Registration Flow}

New users can create an account through the registration form accessible from the login page. The registration process follows these steps:

\begin{enumerate}
    \item Navigate to the application URL and click \textbf{``Create Account''} on the login page.
    \item Fill in the required fields:
    \begin{itemize}
        \item \textbf{Username:} 3--20 characters, alphanumeric and underscores only
        \item \textbf{Email:} Valid email format (used for account recovery)
        \item \textbf{Full Name:} Display name shown in the application
        \item \textbf{Password:} Must meet complexity requirements (see below)
        \item \textbf{Confirm Password:} Must match the password field exactly
    \end{itemize}
    \item Click \textbf{``Register''} to submit the form.
    \item Upon successful registration, the user is redirected to the login page with a confirmation message.
\end{enumerate}

\subsection{Password Requirements}

The application enforces the following password policy:

\begin{table}[htbp]
\centering
\caption{Password complexity requirements}
\label{tab:password-policy}
\begin{tabular}{|l|l|}
\hline
\textbf{Requirement} & \textbf{Details} \\
\hline
Minimum length & 8 characters \\
\hline
Maximum length & 128 characters \\
\hline
Uppercase letters & At least 1 required \\
\hline
Lowercase letters & At least 1 required \\
\hline
Numeric digits & At least 1 required \\
\hline
Special characters & At least 1 required (!@\#\$\%\^{}\&*) \\
\hline
Common passwords & Rejected (checked against blocklist) \\
\hline
\end{tabular}
\end{table}

\noindent A real-time password strength indicator provides visual feedback during registration, displaying strength levels from ``Weak'' (red) through ``Fair'' (orange) to ``Strong'' (green).

\subsection{Login Process}

Returning users authenticate through the login form:

\begin{enumerate}
    \item Enter the registered username or email address.
    \item Enter the account password.
    \item Optionally check \textbf{``Remember me''} to extend session duration.
    \item Click \textbf{``Login''} to authenticate.
\end{enumerate}

\noindent Upon successful authentication, the user is redirected to the Home page and the navigation sidebar becomes fully accessible.

\subsection{Session Persistence}

The application manages user sessions through Streamlit's session state mechanism:

\begin{itemize}
    \item \textbf{Default session:} Expires after 30 minutes of inactivity.
    \item \textbf{``Remember me'' session:} Persists for 7 days via a browser cookie.
    \item \textbf{Manual logout:} Users can explicitly end their session via the sidebar logout button.
    \item \textbf{Session tokens:} Cryptographically generated tokens are stored server-side and validated on each request.
\end{itemize}

\subsection{Account Lockout Mechanism}

To prevent brute-force attacks, the application implements progressive account lockout:

\begin{enumerate}
    \item After \textbf{3 failed attempts}: A CAPTCHA challenge is displayed.
    \item After \textbf{5 failed attempts}: The account is locked for 15 minutes.
    \item After \textbf{10 failed attempts}: The account is locked for 60 minutes and an email notification is sent (if configured).
\end{enumerate}

\noindent The lockout counter resets after a successful login. Administrators can manually unlock accounts through the database if needed:


\subsection{Password Reset}

Users who forget their password can initiate a reset through the \textbf{``Forgot Password?''} link on the login page. The reset flow generates a time-limited token (valid for 1 hour) and displays it to the user for demonstration purposes. In a production environment, this token would be sent via email.

\section{Home Page}
\label{sec:user-home}

The Home page serves as the application's landing page and adapts its content based on the user's authentication status.

\subsection{Before Login (Public View)}

Unauthenticated visitors see:

\begin{itemize}
    \item \textbf{Application title and tagline:} ``Urban Air Quality Forecasting --- Monitor, Analyze, and Predict Air Pollution''
    \item \textbf{Feature highlights:} A grid of cards summarizing key capabilities (data exploration, ML predictions, real-time dashboard, history tracking)
    \item \textbf{Quick statistics:} Aggregate metrics such as total predictions made, active users, and supported pollutants
    \item \textbf{Call-to-action buttons:} Prominent ``Login'' and ``Register'' buttons directing users to the authentication page
    \item \textbf{Technology stack:} Brief overview of the technologies powering the application
\end{itemize}

\subsection{After Login (Authenticated View)}

Authenticated users see a personalized dashboard summary:

\begin{itemize}
    \item \textbf{Welcome message:} Personalized greeting with the user's display name and last login timestamp
    \item \textbf{Activity summary cards:} Number of predictions made, datasets uploaded, and days since registration
    \item \textbf{Recent predictions:} A compact table showing the last 5 predictions with timestamps and model types
    \item \textbf{Quick navigation:} Large icon buttons providing one-click access to each application module
    \item \textbf{Air quality tip of the day:} Rotating educational content about air pollution and health impacts
    \item \textbf{System announcements:} Any active notifications from administrators (e.g., scheduled maintenance)
\end{itemize}

\noindent The Home page layout uses Streamlit's column system to create a responsive grid that adapts to different screen sizes. Key metrics are displayed using \texttt{st.metric()} components with delta indicators showing changes from the previous period.


\section{Data Explorer}
\label{sec:user-explorer}

The Data Explorer module provides tools for uploading, visualizing, and statistically analyzing air quality datasets. It is designed to help users understand their data before proceeding to the prediction module.

\subsection{Data Upload Process}

Users can upload datasets through the file uploader widget:

\begin{enumerate}
    \item Navigate to the \textbf{``Data Explorer''} page from the sidebar.
    \item Click the \textbf{``Browse files''} button or drag and drop a file into the upload area.
    \item The application validates the file format and displays a preview of the first 10 rows.
    \item If validation succeeds, the full dataset is loaded into memory for analysis.
\end{enumerate}

\subsection{Supported File Formats}

\begin{table}[htbp]
\centering
\caption{Supported data file formats}
\label{tab:file-formats}
\begin{tabular}{|l|l|l|l|}
\hline
\textbf{Format} & \textbf{Extension} & \textbf{Max Size} & \textbf{Notes} \\
\hline
CSV & .csv & 50 MB & Comma, semicolon, or tab delimited \\
\hline
Excel & .xlsx, .xls & 50 MB & First sheet read by default \\
\hline
JSON & .json & 50 MB & Records or columnar orientation \\
\hline
Parquet & .parquet & 100 MB & Recommended for large datasets \\
\hline
\end{tabular}
\end{table}

\noindent The application automatically detects the delimiter for CSV files and handles common encoding issues (UTF-8, Latin-1). Column names are normalized by stripping whitespace and converting to lowercase for consistency.

\subsection{Correlation Heatmap}

The correlation heatmap visualizes pairwise Pearson correlation coefficients between all numeric columns in the dataset:

\begin{itemize}
    \item \textbf{Color scale:} Diverging color map from blue ($-1$) through white ($0$) to red ($+1$).
    \item \textbf{Annotations:} Each cell displays the correlation coefficient rounded to two decimal places.
    \item \textbf{Interactivity:} Hovering over a cell shows the exact value and the two correlated variables.
    \item \textbf{Filtering:} Users can select a subset of columns to include in the heatmap via a multiselect widget.
    \item \textbf{Threshold highlighting:} An optional slider allows users to highlight only correlations above a specified absolute threshold.
\end{itemize}

\noindent The heatmap is rendered using Plotly, enabling zoom, pan, and export to PNG. This visualization is particularly useful for identifying multicollinearity between pollutant measurements and meteorological variables.

\subsection{Histogram Analysis}

The histogram module allows users to explore the distribution of individual variables:

\begin{itemize}
    \item \textbf{Column selection:} A dropdown menu lists all numeric columns available for plotting.
    \item \textbf{Bin count:} Adjustable slider (10--100 bins) to control histogram granularity.
    \item \textbf{Distribution overlay:} Optional kernel density estimation (KDE) curve overlaid on the histogram.
    \item \textbf{Normalization:} Toggle between raw counts and probability density.
    \item \textbf{Group by:} Optional categorical column for creating grouped/stacked histograms.
    \item \textbf{Log scale:} Toggle for logarithmic y-axis, useful for skewed distributions.
\end{itemize}

\noindent Users can generate histograms for multiple columns simultaneously using the ``Plot All'' button, which creates a grid of subplots for quick visual comparison.

\subsection{Statistical Summary}

The statistical summary panel provides comprehensive descriptive statistics:

\begin{table}[htbp]
\centering
\caption{Statistical metrics computed for each numeric column}
\label{tab:stats-summary}
\begin{tabular}{|l|l|}
\hline
\textbf{Metric} & \textbf{Description} \\
\hline
Count & Number of non-null values \\
\hline
Mean & Arithmetic mean \\
\hline
Std & Standard deviation \\
\hline
Min / Max & Minimum and maximum values \\
\hline
25\% / 50\% / 75\% & Quartile values \\
\hline
Skewness & Measure of distribution asymmetry \\
\hline
Kurtosis & Measure of distribution tail heaviness \\
\hline
Missing (\%) & Percentage of null/NaN values \\
\hline
\end{tabular}
\end{table}

\noindent The summary is displayed as an interactive dataframe with sortable columns. Users can download the statistical summary as a CSV file using the \textbf{``Download Summary''} button.

\subsection{Additional Explorer Features}

\begin{itemize}
    \item \textbf{Data type detection:} Automatic identification of datetime, numeric, and categorical columns with suggested type conversions.
    \item \textbf{Missing value visualization:} A bar chart showing the percentage of missing values per column, with options to view the missingness pattern matrix.
    \item \textbf{Scatter plot matrix:} Pairwise scatter plots for selected variables with optional color coding by a categorical variable.
    \item \textbf{Time series plot:} If a datetime column is detected, users can plot any numeric variable over time with resampling options (hourly, daily, weekly, monthly).
    \item \textbf{Outlier detection:} IQR-based outlier flagging with visual indicators on box plots.
\end{itemize}


\section{Predict Future Quality}
\label{sec:user-predict}

The Predict Future Quality module is the core forecasting engine of the application. It allows users to upload new air quality measurements and generate predictions using pre-trained machine learning models.

\subsection{Upload Schema Requirements}

The prediction module requires input data to conform to a specific schema. The uploaded CSV file must contain the following columns:

\begin{table}[htbp]
\centering
\caption{Required columns for prediction input}
\label{tab:predict-schema}
\begin{tabular}{|l|l|l|l|}
\hline
\textbf{Column Name} & \textbf{Type} & \textbf{Unit} & \textbf{Required} \\
\hline
datetime & datetime & ISO 8601 & Yes \\
\hline
pm25 & float & $\mu g/m^3$ & Yes \\
\hline
pm10 & float & $\mu g/m^3$ & Yes \\
\hline
no2 & float & ppb & Yes \\
\hline
so2 & float & ppb & Conditional \\
\hline
co & float & ppm & Conditional \\
\hline
o3 & float & ppb & Conditional \\
\hline
temperature & float & \textdegree C & Recommended \\
\hline
humidity & float & \% & Recommended \\
\hline
wind\_speed & float & m/s & Recommended \\
\hline
\end{tabular}
\end{table}

\noindent Columns marked ``Conditional'' are required depending on the selected prediction target. Columns marked ``Recommended'' improve prediction accuracy but are not strictly required---the model will use default imputation strategies for missing meteorological features.

\subsection{Schema Validation}

Upon upload, the application performs comprehensive validation:

\begin{enumerate}
    \item \textbf{Column presence check:} Verifies all required columns exist (case-insensitive matching).
    \item \textbf{Data type validation:} Ensures numeric columns contain parseable values.
    \item \textbf{Range validation:} Flags physically impossible values (e.g., negative concentrations, humidity $> 100\%$).
    \item \textbf{Temporal ordering:} Checks that datetime values are monotonically increasing.
    \item \textbf{Minimum row count:} Requires at least 24 rows (representing 24 hours of data) for meaningful predictions.
\end{enumerate}

\noindent Validation errors are displayed as an itemized list with specific row/column references, enabling users to quickly identify and correct issues in their data.

\subsection{Model Selection}
